\section{KVM}
\label{sec:kvm}

This section presents a comprehensive analysis of execution latencies,
determinism, and real-time capabilities within a KVM-based virtualized
environment. We evaluate system behavior across four host--guest kernel
combinations, ranging from standard non-real-time configurations to
fully optimized end-to-end \texttt{PREEMPT\_RT} stacks, and extend the
analysis to TACLe real-time benchmark workloads under noisy-neighbour
conditions.


\subsection{BASELINE PERFORMANCE ANALYSIS IN KVM
ENVIRONMENTS}\label{baseline-performance-analysis-in-kvm-environments}

This section presents a detailed analysis of execution latencies
measured within a virtualized environment based on the Kernel-based
Virtual Machine (KVM) hypervisor. The primary objective is to establish
a performance baseline and evaluate the system's behavior, determinism,
and virtualization overhead under standard scheduling policies before
introducing more advanced or restrictive tuning configurations.

To quantify response times, internal jitter, and the Worst-Case
Execution Time (WCET), the \texttt{cyclictest} tool was employed through
continuous 5-minute baseline executions. The investigation explores the
impact on system stability and temporal predictability by
cross-referencing four distinct kernel combinations between the
underlying Host infrastructure and the Guest virtual machine:

\begin{itemize}
\item
  \textbf{Non-Real-Time (NRT) Host and Non-Real-Time (NRT) Guest:} to
  measure standard system behavior and baseline virtualization overhead
  in the total absence of real-time optimizations.
\item
  \textbf{Non-Real-Time (NRT) Host and Real-Time (RT) Guest:} to
  evaluate the effectiveness of internal scheduling optimizations within
  the guest when the underlying hypervisor lacks deterministic
  guarantees.
\item
  \textbf{Real-Time (RT) Host and Non-Real-Time (NRT) Guest:} to analyze
  the impact of a determinism-optimized host on the performance,
  preemption spikes, and overall execution stability of a standard,
  general-purpose guest.
\item
  \textbf{Real-Time (RT) Host and Real-Time (RT) Guest (Full RT Stack):}
  to observe the maximum level of temporal predictability achievable by
  aligning both the host infrastructure and the guest operating system
  with the \texttt{PREEMPT\_RT} patch.
\end{itemize}

The analyses in the following paragraphs offer a comprehensive overview
of the latency profiles characteristic of each configuration, detailing
the limitations of general-purpose environments and validating the
effectiveness of an end-to-end real-time stack for mitigating
virtualization overhead.

\begin{figure}
\centering
\includesvg[width=\linewidth]{../KVM/tests/plot/svg/kvm_nonoise.svg}
\caption{Baseline Performance - KVM}
\end{figure}

\subsubsection{NON-REAL-TIME HOST AND NON-REAL-TIME
GUEST}\label{non-real-time-host-and-non-real-time-guest}

This section presents an analysis of a single, extended 5-minute
\texttt{cyclictest} run conducted in a standard environment, featuring a
Non-Real-Time guest Linux kernel hosted on a Non-Real-Time host system.
The objective is to establish a performance baseline and assess system
determinism in the absence of real-time optimizations, such as the
\texttt{PREEMPT\_RT} patch or a real-time tuned host hypervisor.

\paragraph{Distribution Analysis and Nominal
Performance}\label{distribution-analysis-and-nominal-performance}

Analysis of the histogram log reveals significant consistency during
nominal execution over the 5-minute test period. The statistical mode
(the most frequent latency value) was anchored firmly at 6 $\mu s$. The
average latency was equally stable, measuring exactly 7 $\mu s$ across the
entire continuous execution. Additionally, the minimum recorded latency
was extremely low, bottoming out at 5 $\mu s$. This performance indicates
that under nominal conditions, without background load or anomalous
preemption, the baseline overhead introduced by the default
virtualization layer and host kernel scheduler is remarkably low.

\paragraph{Worst-Case Execution Time (WCET) and Lack of
Determinism}\label{worst-case-execution-time-wcet-and-lack-of-determinism}

The fundamental limitation of this general-purpose configuration is
revealed by the Worst-Case Execution Time (WCET) data, which continues
to demonstrate high variability and unpredictable latency spikes. Across
the 5-minute continuous execution, the absolute maximum latency recorded
was 602 $\mu s$, a severe, stochastic deviation from the average.

\paragraph{Conclusions on Standard
Environments}\label{conclusions-on-standard-environments}

These latency spikes---the ``long tails'' observed in the distribution
logs---are characteristic of general-purpose software stacks. In this
ecosystem, the hypervisor operates without strict real-time constraints
and may preempt the Virtual CPU (VCPU) to serve host-level tasks;
simultaneously, the guest kernel is subject to non-preemptible critical
sections and non-deferrable hardware interrupts. Consequently, while
nominal and average performance metrics are excellent, the Non-RT guest
on a Non-RT host environment is susceptible to unpredictable delays and
cannot guarantee the rigid, reliable upper bounds required for
safety-critical control applications.

\subsubsection{NON-REAL-TIME HOST AND REAL-TIME
GUEST}\label{non-real-time-host-and-real-time-guest}

This section analyzes the results obtained from a single, prolonged
5-minute execution of the \texttt{cyclictest} tool in a hybrid
environment, configured with a Real-Time optimized guest hosted on a
Standard Linux Kernel Host. The goal of this analysis is to evaluate if
the guest's internal real-time scheduling can maintain determinism
during sustained workloads when the underlying host hypervisor lacks
real-time optimizations.

\paragraph{Nominal Performance and Average
Latency}\label{nominal-performance-and-average-latency}

The data collected confirm that optimization at the guest level
guarantees excellent efficiency in the average case. The average latency
value is extremely stable, resting at 8 $\mu s$ for the duration of the
5-minute test. Furthermore, the absolute minimum latency dropped down to
4 $\mu s$. This behavior indicates that the host allocates VCPU resources to
the guest promptly and with minimal virtualization overhead under normal
operating conditions despite not being patched.

\paragraph{Worst-Case Execution Time (WCET)
Analysis}\label{worst-case-execution-time-wcet-analysis}

Despite the high baseline efficiency of the underlying infrastructure,
the analysis of the Worst-Case Execution Time (WCET) reveals a severe
lack of determinism under sustained testing. The recorded maximum
latency reached an extreme peak of 6114 $\mu s$ (over 6 milliseconds), and
other 3 over-milliseconds spikes. This data indicates that while the
system performs well nominally, it is subject to rare but catastrophic
scheduling delays that shatter any real-time guarantees.

\paragraph{Conclusions on the Hybrid
Environment}\label{conclusions-on-the-hybrid-environment}

The empirical observation of this sustained 5-minute test demonstrates
conclusively that the exclusive optimization of the virtualization
infrastructure (Hypervisor/Host) is not able to stem anomalous latencies
if the guest operating system is not equally optimized. The critical
peak of over 6 milliseconds is a direct symptom of the internal
architecture of the standard kernel. The absence of the
\texttt{PREEMPT\_RT} patch leaves the guest vulnerable to delays induced
by the host's non-preemptible critical sections (such as spinlocks),
non-deferrable hardware interrupts, and unbounded priority inversion
phenomena.

Therefore, to obtain reliable determinism in control contexts, the use
of an RT Guest proves not to be a sufficient condition.

\subsubsection{REAL-TIME HOST AND NON-REAL-TIME
GUEST}\label{real-time-host-and-non-real-time-guest}

This section analyzes the results obtained from a single, prolonged
5-minute execution of the \texttt{cyclictest} tool in a hybrid
environment, configured with a Standard Linux Guest hosted on a
Real-Time \texttt{PREEMPT\_RT} patched Host. The goal of this analysis
is to verify whether the use of a deterministic Host hypervisor is
sufficient to guarantee strict temporal constraints during sustained
workloads when the guest operating system remains general-purpose.

\paragraph{Nominal Performance and Average
Latency}\label{nominal-performance-and-average-latency-1}

The data collected confirm that optimization at the Guest OS level
guarantees excellent efficiency in the average case. The average latency
value is extremely stable, resting at 8 $\mu s$ for the duration of the
5-minute test. Furthermore, the absolute minimum latency dropped down to
4 $\mu s$. This behavior indicates that the Real-Time guest handles thread
scheduling promptly and with minimal overhead under normal operating
conditions.

\paragraph{Worst-Case Execution Time (WCET)
Analysis}\label{worst-case-execution-time-wcet-analysis-1}

The analysis of the Worst-Case Execution Time (WCET) reveals highly
stable and deterministic behavior during this sustained test. The
absolute maximum latency recorded reached a peak of only 93 $\mu s$. This
data demonstrates a complete absence of the severe, prolonged
hypervisor-induced preemption spikes that typically characterize
standard host environments.

\paragraph{Conclusions on the Hybrid
Environment}\label{conclusions-on-the-hybrid-environment-1}

The empirical observation of this sustained 5-minute test demonstrates
conclusively that the exclusive optimization of the guest operating
system is not able to stem anomalous latencies if the virtualization
infrastructure (Hypervisor/Host) is not equally optimized. The critical
peak of over 6 milliseconds is a direct symptom of the internal
architecture of the standard host kernel. The absence of the
\texttt{PREEMPT\_RT} patch on the host leaves the guest vulnerable to
delays induced by the host's non-preemptible critical sections (such as
spinlocks), non-deferrable hardware interrupts, and unbounded priority
inversion phenomena.

\subsubsection{REAL-TIME HOST AND REAL-TIME GUEST (FULL RT
STACK)}\label{real-time-host-and-real-time-guest-full-rt-stack}

This section analyzes the results obtained from a single, prolonged
5-minute execution of the \texttt{cyclictest} tool in a ``Full RT''
environment, where a \texttt{PREEMPT\_RT} patched Linux guest kernel is
hosted on a Real-Time optimized Host. The objective of this
configuration is to validate the effectiveness of an end-to-end
real-time stack in mitigating virtualization overhead and establishing a
strict, deterministic upper bound for execution latency during a
sustained workload.

\paragraph{Nominal Performance and Average
Latency}\label{nominal-performance-and-average-latency-2}

The experimental data collected demonstrates unparalleled stability
under sustained nominal conditions. The average latency metric is
exceptionally consistent, locked at 8 $\mu s$ for the duration of the test.
The absolute minimum latency recorded was 7 $\mu s$. The histogram
distribution highlights this extreme efficiency, showing that the vast
majority of execution cycles---over 92\% of samples---completed in
exactly 8 $\mu s$. This indicates that when the Virtual CPU (VCPU) is active,
the \texttt{PREEMPT\_RT} patch successfully manages thread scheduling
and interrupt handling with virtually zero internal jitter, effectively
mirroring bare-metal performance for standard task cycles.

\paragraph{Worst-Case Execution Time (WCET)
Analysis}\label{worst-case-execution-time-wcet-analysis-2}

The transition to a Full RT stack demonstrates a definitive improvement
in worst-case performance, completely neutralizing the severe stochastic
latency spikes observed in non-optimized configurations. Throughout the
entire 5-minute sustained test, the absolute maximum latency (Worst-Case
Execution Time) was strictly bounded at just 81 $\mu s$. The ``long tail'' of
the distribution---characteristic of hypervisor-induced preemption and
unoptimized kernel locks---has been entirely eliminated. The system
demonstrates a robust, flawless ability to remain within deterministic
bounds, drastically reducing the magnitude of any scheduling delays.

\paragraph{Conclusions on the Full RT
Architecture}\label{conclusions-on-the-full-rt-architecture}

The empirical observation of this sustained test provides conclusive
evidence regarding the requirements for deterministic virtualization:

\begin{itemize}
\item
  \textbf{End-to-End Determinism:} Achieving hard real-time guarantees
  is an end-to-end property, and the combination of the
  \texttt{PREEMPT\_RT} Guest kernel and a Real-Time Host successfully
  prevents hypervisor-induced preemption.
\item
  \textbf{Absolute Stability:} The system maintained a strict upper
  bound of 81 $\mu s$ over nearly 6 million consecutive execution cycles.
\item
  \textbf{Zero Overflows:} The complete absence of histogram overflows
  confirms that the system never experienced uncontrolled latency spikes
  during the prolonged execution.
\item
  \textbf{Viability for Critical Systems:} This ``Full RT''
  configuration establishes a highly predictable WCET, proving that with
  proper infrastructure tuning, virtualized environments can reliably
  support safety-critical control applications that previously required
  dedicated bare-metal hardware.
\end{itemize}

\begin{figure}
\centering
\includesvg[width=\linewidth]{../KVM/tests/plot/svg/KVM_baseline_boxplot.svg}
\caption{Baseline Performance - KVM}
\end{figure}

\subsection{IMPACT OF THE STRESS WORKLOAD ON
LATENCIES}\label{impact-of-the-stress-workload-on-latencies}

To evaluate system robustness and trigger potentially higher latencies,
the testing methodology involves introducing an additional load, defined
as a ``stress workload.'' In the case of the KVM hypervisor, this stress
workload is executed in the background directly on the Host operating
system, simulating heavy external contention that competes for CPU time
and system resources.

Experimental analysis of the KVM infrastructure reveals that adding this
background load produces severe performance degradation across the
board. Unlike the complex, configuration-dependent behaviors observed in
other hypervisors, the stress workload on KVM exposes fundamental
architectural vulnerabilities in maintaining determinism:

\begin{itemize}
\item
  \textbf{Catastrophic Failure in Unoptimized and Hybrid Stacks:}
  Introducing host-level stress immediately shatters the determinism of
  standard and hybrid configurations. Whether using a fully
  Non-Real-Time (NRT) stack or attempting partial optimization (an RT
  Guest on an NRT Host, or an NRT Guest on an RT Host), the system
  experiences catastrophic scheduling delays. Maximum latencies in these
  configurations routinely spike between 44 and 51 milliseconds,
  indicating massive host-induced starvation and lock contention.
\item
  \textbf{The Illusion of Guest-Only Optimization:} Counter-intuitively,
  equipping the Guest with a Real-Time kernel while the Host remains NRT
  yielded the highest variability and the worst absolute latency spikes
  (over 51 milliseconds) of the dataset. This proves that an optimized
  guest scheduler is entirely powerless to protect time-sensitive
  threads if the underlying Host hypervisor is vulnerable to
  non-deterministic preemption.
\item
  \textbf{Chronic Instability of the Full RT Stack:} Even when utilizing
  an end-to-end ``Full RT'' stack (RT Host and RT Guest) with optimal
  prioritization parameters, the KVM environment fails to guarantee
  strict real-time bounds under stress. While the Full RT stack
  successfully suppresses the catastrophic 50-millisecond delays
  (bounding the absolute maximum latency to approximately 9.8
  milliseconds), it suffers from a massive frequency of significant
  scheduling stalls, evidenced by hundreds of over-millisecond latency
  overflows.
\end{itemize}

In this section, we will present the detailed experiments conducted
under these stress conditions and analyze the resulting execution logs.

\begin{figure}
\centering
\includesvg[width=\linewidth]{../KVM/tests/plot/svg/kvm_backgroundnoise.svg}
\caption{Performance under stress workload- KVM at maximum priority}
\end{figure}

\subsubsection{NON REAL-TIME HOST AND NON REAL-TIME GUEST -STRESS
HOST}\label{non-real-time-host-and-non-real-time-guest--stress-host}

This section presents an analysis of the first \texttt{cyclictest}
execution log conducted on a KVM host under stress conditions. The
objective is to evaluate the system's scheduling behavior and latency
bounds when subjected to external load.

\paragraph{Distribution Analysis and Nominal
Performance}\label{distribution-analysis-and-nominal-performance-1}

Analysis of the histogram log reveals that the system maintains a
reasonable baseline under stress. Over the test duration, the average
latency recorded was 13 $\mu s$. The statistical mode (the most frequent
latency) was concentrated at 14 $\mu s$, with a significant number of cycles
also completing at 13 $\mu s$. The absolute minimum recorded latency dropped
to an impressive 5 $\mu s$. This indicates that when the CPU is not actively
dealing with stress-induced contention, the base virtualization and
scheduling overhead remains low.

\paragraph{Worst-Case Execution Time (WCET) and Lack of
Determinism}\label{worst-case-execution-time-wcet-and-lack-of-determinism-1}

The impact of the stress workload becomes severely apparent when
analyzing the Worst-Case Execution Time (WCET). The absolute maximum
latency recorded during this run was an extreme 44,409 $\mu s$ (over 44
milliseconds). Furthermore, the log reports 11 histogram overflows
(latencies exceeding the 1000 $\mu s$ tracking threshold).

\paragraph{Conclusions}\label{conclusions}

The presence of 44-millisecond latency spikes clearly demonstrates that
this environment lacks determinism. Under stress, the host OS
experiences massive scheduling delays, likely due to lock contention,
non-preemptible critical sections, or starvation of the test threads.
This configuration cannot guarantee the strict upper bounds required for
real-time applications.

\subsubsection{NON REAL-TIME HOST AND REAL-TIME GUEST -STRESS
HOST}\label{non-real-time-host-and-real-time-guest--stress-host}

This section analyzes the results obtained from the second
\texttt{cyclictest} execution, also run under stressed conditions on the
KVM host infrastructure.

\paragraph{Nominal Performance and Average
Latency}\label{nominal-performance-and-average-latency-3}

The nominal performance mirrors the first log closely, confirming a
consistent baseline efficiency even under load. The average latency was
precisely 13 $\mu s$, and the minimum latency was 5 $\mu s$. Interestingly, the
distribution shows a dual-peak behavior, with the highest concentration
at 13 $\mu s$ (the mode), but with another massive spike of cycles completing
at 7 $\mu s$.

\paragraph{Worst-Case Execution Time (WCET)
Analysis}\label{worst-case-execution-time-wcet-analysis-3}

This specific execution experienced the highest variability and worst
latency spikes of the entire dataset. The absolute maximum latency
reached an exorbitant 51,069 $\mu s$ (over 51 milliseconds). The severity of
the instability is further highlighted by the 26 histogram overflows,
more than double the number seen in the first log.

\paragraph{Conclusions}\label{conclusions-1}

The empirical observation of this run conclusively shows that the stress
workload induces chaotic, non-deterministic preemption. A 51-millisecond
delay represents a catastrophic failure for any safety-critical system.
The KVM host in this state is entirely unsuitable for bounded real-time
execution, as the scheduler cannot adequately protect time-sensitive
threads from background stress.

\subsubsection{REAL-TIME HOST AND NON REAL-TIME GUEST -STRESS
HOST}\label{real-time-host-and-non-real-time-guest--stress-host}

This section examines the third prolonged execution of the
\texttt{cyclictest} tool on the stressed KVM host environment,
validating the patterns observed in previous iterations.

\paragraph{Nominal Performance and Average
Latency}\label{nominal-performance-and-average-latency-4}

The data collected confirms a stable nominal execution path. The average
latency metric remained locked at 13 $\mu s$, while the minimum latency
reached the lowest point among the first three runs, hitting 3 $\mu s$. The
histogram shows a massive concentration of execution cycles completing
in exactly 13 $\mu s$, indicating that when the scheduler is unimpeded, it
processes tasks with high consistency.

\paragraph{Worst-Case Execution Time (WCET)
Analysis}\label{worst-case-execution-time-wcet-analysis-4}

Despite the strong nominal performance, the WCET remains unacceptable
for real-time standards. The maximum latency peaked at 50,257 $\mu s$ (over
50 milliseconds). The system registered 11 histogram overflows,
indicating multiple occurrences of severe scheduling stalls exceeding
one millisecond.

\paragraph{Conclusions}\label{conclusions-2}

This run confirms that while the system can occasionally achieve
extremely fast task dispatching (3 $\mu s$ min), it is completely vulnerable
to unpredictable, systemic delays. The 50-millisecond peak re-emphasizes
that background stress can effectively starve tasks for dozens of
milliseconds, a fatal condition for latency-sensitive applications.

\subsubsection{REAL-TIME HOST AND REAL-TIME GUEST -STRESS
HOST}\label{real-time-host-and-real-time-guest--stress-host}

This section presents an analysis of the provided \texttt{cyclictest}
execution log conducted on a KVM environment under stress conditions.
The objective is to evaluate the system's scheduling determinism and
latency bounds when subjected to external load.

\paragraph{Distribution Analysis and Nominal
Performance}\label{distribution-analysis-and-nominal-performance-2}

Analysis of the histogram log reveals a highly efficient baseline
performance under nominal execution. The statistical mode (the most
frequent latency) is sharply concentrated at 7 $\mu s$, representing over 1.7
million completed cycles. The average latency metric is remarkably
stable at 11 $\mu s$, and the absolute minimum recorded latency is
exceptionally low at 4 $\mu s$. This indicates that most of the time, the
virtualization overhead is minimal.

\paragraph{Worst-Case Execution Time (WCET)
Analysis}\label{worst-case-execution-time-wcet-analysis-5}

Despite the excellent average case, the analysis of the Worst-Case
Execution Time (WCET) reveals severe scheduling instability induced by
the stress workload. The absolute maximum latency recorded reached 9,829
$\mu s$ (nearly 9.8 milliseconds). Even more critically, the log reports an
overwhelming 283 histogram overflows (latencies exceeding the 1000 $\mu s$
tracking threshold). This represents a massive frequency of significant
scheduling stalls compared to typical runs.

\paragraph{Conclusions}\label{conclusions-3}

The empirical observation of this execution demonstrates severe latency
instability, which is particularly critical given the system's strict
configuration. Despite this optimal prioritization, the system still
suffered from chronic and significant delays.

These high latencies indicate that the delays are originating from
deeper, non-preemptible sources escaping the guest's control.
Consequently, this configuration proves that merely applying the maximum
scheduler priority is fundamentally insufficient to guarantee the
strict, reliable upper bounds required for safety-critical real-time
applications.

\begin{figure}
\centering
\includesvg[width=\linewidth]{../KVM/tests/plot/svg/KVM_Stress_Workload_boxplot.svg}
\caption{Performance under stress workload- KVM at maximum priority}
\end{figure}

\subsection{ISOLATED STRESS HOST PERFORMANCE
ANALYSIS}\label{isolated-stress-host-performance-analysis}

This section presents a detailed analysis of execution latencies
measured within a virtualized environment based on the Kernel-based
Virtual Machine (KVM) hypervisor, subjected to a background stress
workload but utilizing OS-level isolation strategies. The primary
objective is to evaluate whether spatial isolation of the Virtual CPUs
(VCPUs) can effectively shield the guest operating system from the
severe preemption and scheduling delays induced by host-level
contention.

To quantify response times, internal jitter, and the Worst-Case
Execution Time (WCET) under these isolated conditions, the
\texttt{cyclictest} tool was employed through continuous 5-minute
executions. The investigation explores the impact on system stability
and temporal predictability by cross-referencing four distinct kernel
combinations between the underlying Host infrastructure and the Guest
virtual machine:

\begin{itemize}
\item
  \textbf{Non-Real-Time (NRT) Host and Non-Real-Time (NRT) Guest:} to
  measure system behavior and virtualization overhead when applying
  isolation techniques to a standard, unoptimized software stack.
\item
  \textbf{Non-Real-Time (NRT) Host and Real-Time (RT) Guest:} to
  evaluate the effectiveness of an optimized guest scheduler operating
  on dedicated cores, while the underlying hypervisor lacks
  deterministic guarantees.
\item
  \textbf{Real-Time (RT) Host and Non-Real-Time (NRT) Guest:} to analyze
  the impact of a determinism-optimized host on the preemption spikes of
  a standard guest when spatial isolation is enforced.
\item
  \textbf{Real-Time (RT) Host and Real-Time (RT) Guest (Full RT Stack):}
  to observe the maximum level of temporal predictability achievable by
  combining CPU isolation with an end-to-end \texttt{PREEMPT\_RT}
  patched infrastructure.
\end{itemize}

The analyses in the following paragraphs offer a comprehensive overview
of the latency profiles characteristic of each isolated configuration,
detailing the effectiveness of core pinning in mitigating the
catastrophic failures observed in un-isolated stress scenarios.

\begin{figure}
\centering
\includesvg[width=\linewidth]{../KVM/tests/plot/svg/kvm_backgroundnoise_isolated.svg}
\caption{Performance under stress- KVM OS-level isolation}
\end{figure}

\subsubsection{NON-REAL-TIME HOST AND NON-REAL-TIME GUEST (ISOLATED
STRESS
HOST)}\label{non-real-time-host-and-non-real-time-guest-isolated-stress-host}

This section presents an analysis of a single, extended 5-minute
\texttt{cyclictest} run conducted in an isolated environment, featuring
a Non-Real-Time guest Linux kernel hosted on a Non-Real-Time host system
under stress. The objective is to establish a performance baseline for
spatial isolation without real-time kernel patches.

\paragraph{Distribution Analysis and Nominal
Performance}\label{distribution-analysis-and-nominal-performance-3}

Analysis of the histogram log reveals a very consistent nominal
execution under isolated stress. The statistical mode (the most frequent
latency value) was anchored firmly at 8 $\mu s$. The average latency was
highly stable, measuring exactly 11 $\mu s$ across the entire continuous
execution. Additionally, the minimum recorded latency dropped to 6 $\mu s$.
This indicates that CPU isolation successfully protects the active VCPU
during standard execution cycles, keeping the baseline virtualization
overhead low despite the background load on the host.

\paragraph{Worst-Case Execution Time (WCET) and Lack of
Determinism}\label{worst-case-execution-time-wcet-and-lack-of-determinism-2}

While isolation drastically improved the overall stability compared to
un-isolated stress tests (which saw 44-millisecond delays), the
Worst-Case Execution Time (WCET) data reveals that determinism is not
fully guaranteed. The absolute maximum latency recorded was bounded at
2,024 $\mu s$ (roughly 2 milliseconds). However, the system still registered
19 histogram overflows (latencies exceeding the 1000 $\mu s$ threshold).

\paragraph{Conclusions on Standard Isolated
Environments}\label{conclusions-on-standard-isolated-environments}

The application of CPU isolation prevents the catastrophic
multi-millisecond starvation seen in standard stress scenarios, reducing
the WCET by over 95\%. However, the presence of nearly 2-millisecond
spikes and 19 overflows confirms that a general-purpose NRT/NRT stack
remains susceptible to unpredictable delays. Because the tasks cannot be
preempted by standard user-space processes (due to the FIFO=99
scheduling policy), these spikes originate from deeper, non-preemptible
sources such as host kernel lockups, unmaskable hardware interrupts, or
shared hardware resource contention (like cache and memory bus) induced
by the stress load running on neighboring cores.

\subsubsection{NON-REAL-TIME HOST AND REAL-TIME GUEST (ISOLATED STRESS
HOST)}\label{non-real-time-host-and-real-time-guest-isolated-stress-host}

This section analyzes the results obtained from a single, prolonged
5-minute execution of the \texttt{cyclictest} tool in a hybrid
environment, configured with a Real-Time optimized guest hosted on a
Standard Linux Kernel Host, with CPU isolation applied during host
stress.

\paragraph{Nominal Performance and Average
Latency}\label{nominal-performance-and-average-latency-5}

The data collected confirms that the isolated RT guest achieves
excellent baseline efficiency. The average latency value rested at 12
$\mu s$, with the statistical mode concentrated at 9 $\mu s$. Furthermore, the
absolute minimum latency reached an impressive 4 $\mu s$, indicating that the
optimized guest scheduler, when isolated, dispatches tasks with extreme
speed under nominal conditions.

\paragraph{Worst-Case Execution Time (WCET)
Analysis}\label{worst-case-execution-time-wcet-analysis-6}

The analysis of the Worst-Case Execution Time (WCET) reveals a massive
improvement over the un-isolated equivalent (which previously failed
catastrophically with 51-millisecond delays). With isolation, the
absolute maximum latency recorded was capped at 2,299 $\mu s$. However, the
system still experienced 11 histogram overflows over the 1-millisecond
threshold.

\paragraph{Conclusions on the Hybrid Isolated
Environment}\label{conclusions-on-the-hybrid-isolated-environment}

The empirical observation of this test demonstrates that isolating the
VCPUs successfully shields the RT Guest from the chaotic preemption of
the NRT Host's stress workload. Yet, the 2.3-millisecond peak and the 11
overflows indicate a lack of strict determinism. Since the internal
guest tasks are running at maximum real-time priority, the remaining
latency spikes must be attributed to the NRT host infrastructure. The
underlying standard hypervisor still introduces non-deferrable
interrupts or unpredictable virtualization overhead that occasionally
stalls the isolated cores, proving that guest-side optimization alone is
insufficient.

\subsubsection{REAL-TIME HOST AND NON-REAL-TIME GUEST (ISOLATED STRESS
HOST)}\label{real-time-host-and-non-real-time-guest-isolated-stress-host}

This section examines the prolonged execution of the \texttt{cyclictest}
tool in a hybrid environment featuring a Standard Linux Guest hosted on
a Real-Time \texttt{PREEMPT\_RT} patched Host, operating under isolated
stress conditions.

\paragraph{Nominal Performance and Average
Latency}\label{nominal-performance-and-average-latency-6}

The nominal performance metrics indicate a highly stable execution path.
The average latency metric was maintained at 13 $\mu s$, and the statistical
mode was heavily concentrated at 9 $\mu s$. The absolute minimum latency
recorded was 6 $\mu s$. This confirms that the RT Host efficiently manages
the isolated VCPUs, providing a steady execution foundation.

\paragraph{Worst-Case Execution Time (WCET)
Analysis}\label{worst-case-execution-time-wcet-analysis-7}

The Worst-Case Execution Time (WCET) analysis shows the lowest peak
latency among the hybrid configurations. The absolute maximum latency
recorded reached 1,502 $\mu s$ (approximately 1.5 milliseconds). Despite this
tighter upper bound, the system still accumulated 16 histogram overflows
during the 5-minute sustained test.

\paragraph{Conclusions on the Hybrid Isolated
Environment}\label{conclusions-on-the-hybrid-isolated-environment-1}

Optimizing the Host OS with a \texttt{PREEMPT\_RT} kernel, combined with
CPU isolation, yielded a highly robust infrastructure that successfully
limited the maximum delay to 1.5 milliseconds under heavy stress.
However, the 16 recorded overflows demonstrate that the NRT Guest
remains a bottleneck. Unoptimized internal locks and non-preemptible
critical sections within the general-purpose guest kernel still generate
periodic stalls, preventing the system from achieving hard real-time
determinism despite the stable host infrastructure.

\subsubsection{REAL-TIME HOST AND REAL-TIME GUEST (FULL RT STACK -
ISOLATED STRESS
HOST)}\label{real-time-host-and-real-time-guest-full-rt-stack---isolated-stress-host}

This section analyzes the results obtained from the \texttt{cyclictest}
execution in a ``Full RT'' environment, where a \texttt{PREEMPT\_RT}
patched Linux guest kernel is hosted on a Real-Time optimized Host, with
spatial isolation applied during the stress workload. The objective is
to evaluate the absolute limits of temporal predictability in a fully
tuned KVM stack.

\paragraph{Distribution Analysis and Nominal
Performance}\label{distribution-analysis-and-nominal-performance-4}

The experimental data highlights unparalleled efficiency in the
best-case and nominal scenarios. The statistical mode was anchored at 9
$\mu s$, and the average latency was highly stable at 13 $\mu s$. Most notably,
the absolute minimum latency recorded was an exceptional 3 $\mu s$---the
lowest value observed across all configurations. This demonstrates that
when the full RT stack operates unimpeded on isolated cores, the
internal jitter and virtualization overhead are practically negligible.

\paragraph{Worst-Case Execution Time (WCET)
Analysis}\label{worst-case-execution-time-wcet-analysis-8}

Despite the optimal software stack and spatial isolation, the Worst-Case
Execution Time (WCET) analysis reveals that strict hard real-time bounds
are still compromised under stress. The absolute maximum latency was
recorded at 1,935 $\mu s$ (nearly 1.9 milliseconds), and the system
experienced 14 histogram overflows over the 1000 $\mu s$ tracking threshold.

\paragraph{Conclusions on the Full RT Isolated
Architecture}\label{conclusions-on-the-full-rt-isolated-architecture}

The empirical observation of this execution provides critical insight
into the limits of virtualization determinism. Combining a Full RT stack
with CPU isolation successfully mitigated the severe, multi-millisecond
starvation caused by host stress, keeping the WCET under 2 milliseconds.
However, the persistence of 14 overflows and a 1.9-millisecond peak
proves that strict determinism is not perfectly guaranteed.

\begin{figure}
\centering
\includesvg[width=\linewidth]{../KVM/tests/plot/svg/KVM_ISOLATED_STRESS_HOST_boxplot.svg}
\caption{Performance under stress- KVM OS-level isolation}
\end{figure}


Following the detailed analysis of each individual scenario, the table
below provides a consolidated overview of the Worst-Case Execution Time
(WCET) measurements. It allows for a direct comparison across all four
kernel configurations (\textbf{NRT-NRT}, \textbf{NRT-RT},
\textbf{RT-NRT}, and \textbf{RT-RT}) under the three tested conditions:
standard execution (\textbf{BASELINE}), heavy system load in the control
domain (\textbf{STRESS DOM0}), and load with isolation mechanisms
applied (\textbf{STRESS DOM0 ISOLATED}).

\begin{table*}
\centering
\caption{KVM WCET Summary}
\label{tab:kvm:wcet-summary}
\begin{tabularx}{\linewidth}{@{}YYYYY@{}}
\toprule
Configuration & NRT-NRT & NRT-RT & RT-NRT & RT-RT\tabularnewline
\midrule
\textbf{BASELINE} & 602 & 6114 & 93 & 81\tabularnewline
\textbf{STRESSHOST} & 44409 & 51069 & 50257 & 9829\tabularnewline
\textbf{STRESSHOST ISOLATED} & 2024 & 2299 & 1502 & 1935\tabularnewline
\bottomrule
\end{tabularx}
\end{table*}

Furthermore, alongside the WCET table, this section presents a detailed
breakdown of the percentage increments for both the maximum latency
(WCET) and the average latency (expressed as Mean \(\pm\) Standard
Deviation). This provides a precise quantitative analysis of the
performance degradation induced by the stress workload compared to the
baseline for each specific scenario.

\subsection{SUMMARY OF RESULTS}\label{summary-of-results}

\subsubsection{NRT NRT}\label{nrt-nrt}

\begin{table*}
\centering
\caption{NRT-NRT Latency Statistics (KVM)}
\label{tab:kvm:nrt-nrt}
\begin{tabularx}{\linewidth}{@{}YYY@{}}
\toprule
METRIC & BASELINE & STRESS WORKLOAD\tabularnewline
\midrule
\textbf{Average ± SD} & 7.10 ± 2.39 $\mu s$ & 13.60 ± 11.18 $\mu s$\tabularnewline
\textbf{WCET (Max)} & 602 $\mu s$ & 44409 $\mu s$\tabularnewline
\bottomrule
\end{tabularx}
\end{table*}

\textbf{PERCENTAGE INCREMENTS (Stress vs.~Baseline):} * Average
Increment: +91.50\% * WCET Increment: +7276.91\%

\subsubsection{NRT RT}\label{nrt-rt}

\begin{table*}
\centering
\caption{NRT-RT Latency Statistics (KVM)}
\label{tab:kvm:nrt-rt}
\begin{tabularx}{\linewidth}{@{}YYY@{}}
\toprule
METRIC & BASELINE & STRESS WORKLOAD\tabularnewline
\midrule
\textbf{Average ± SD} & 8.55 ± 3.27 $\mu s$ & 13.69 ± 12.00 $\mu s$\tabularnewline
\textbf{WCET (Max)} & 6114 $\mu s$ & 51069 $\mu s$\tabularnewline
\bottomrule
\end{tabularx}
\end{table*}

\textbf{PERCENTAGE INCREMENTS (Stress vs.~Baseline):} * Average
Increment: +60.10\% * WCET Increment: +735.28\%

\subsubsection{RT NRT}\label{rt-nrt}

\begin{table*}
\centering
\caption{RT-NRT Latency Statistics (KVM)}
\label{tab:kvm:rt-nrt}
\begin{tabularx}{\linewidth}{@{}YYY@{}}
\toprule
METRIC & BASELINE & STRESS WORKLOAD\tabularnewline
\midrule
\textbf{Average ± SD} & 7.22 ± 0.98 $\mu s$ & 13.33 ± 8.54 $\mu s$\tabularnewline
\textbf{WCET (Max)} & 93 $\mu s$ & 50257 $\mu s$\tabularnewline
\bottomrule
\end{tabularx}
\end{table*}

\textbf{PERCENTAGE INCREMENTS (Stress vs.~Baseline):} * Average
Increment: +84.76\% * WCET Increment: +53939.78\%

\subsubsection{RT RT}\label{rt-rt}

\begin{table*}
\centering
\caption{RT-RT Latency Statistics (KVM)}
\label{tab:kvm:rt-rt}
\begin{tabularx}{\linewidth}{@{}YYY@{}}
\toprule
METRIC & BASELINE & STRESS WORKLOAD\tabularnewline
\midrule
\textbf{Average ± SD} & 8.15 ± 1.15 $\mu s$ & 11.84 ± 12.79 $\mu s$\tabularnewline
\textbf{WCET (Max)} & 81 $\mu s$ & 9829 $\mu s$\tabularnewline
\bottomrule
\end{tabularx}
\end{table*}

\textbf{PERCENTAGE INCREMENTS (Stress vs.~Baseline):} * Average
Increment: +45.20\% * WCET Increment: +12034.57\%

\subsection{TACLe BENCHMARK}\label{tacle-benchmark-2}

This section of the documentation illustrates the rationale behind the
selection of the benchmarks extracted from the \textbf{TACLeBench}
version 1.9 suite and the methodology adopted for their execution within
our architecture. The goal is to provide a heterogeneous workload to
accurately validate execution latencies and performance stability in
environments with strict real-time requirements.

\subsubsection{1. Benchmark Selection}\label{benchmark-selection}

We selected a representative program for each of the main TACLeBench
categories, ensuring optimal coverage of the different execution
patterns.

\begin{itemize}
\item
  \textbf{Kernel Benchmark (\texttt{matrix1}):} Isolates the most
  computationally intensive portions of code to evaluate raw CPU
  performance and cache efficiency.
\item
  \textbf{Sequential Benchmark (\texttt{huff\_enc}):} Evaluates
  sequential processing and memory access patterns. Data compression
  (325 SLOC, David Bourgin) is excellent for measuring latency
  variations in single-threaded execution.
\item
  \textbf{Test Benchmark (\texttt{test3}):} An artificial stress test
  for WCET (Worst-Case Execution Time, 4235 SLOC, Universität des
  Saarlandes) analysis. It pushes the execution engine to its limits to
  measure time safety margins and system robustness.
\item
  \textbf{Parallel Benchmark (\texttt{Debie}):} Aerospace observation
  tool (6615 SLOC, Tidorum Ltd) consisting of 8 tasks. Essential for
  testing synchronization mechanisms and preemption in multi-tasking
  scenarios.
\item
  \textbf{Application Benchmark (\texttt{lift}):} An elevator controller
  (361 SLOC, Martin Schoeberl). Verifies that the latency metrics from
  synthetic tests guarantee stability in a real cyber-physical control
  application.
\end{itemize}

\subsubsection{2. Execution Methodology and Test
Scenarios}\label{execution-methodology-and-test-scenarios}

To rigorously analyze the system's behavior and the impact of the
virtualization architecture, \textbf{all 5 selected benchmarks were
executed on the RT-RT (Real Time Host, Real-Time Guest) configuration}.

For each benchmark, we defined a test matrix composed of four
operational scenarios. In all configurations, \textbf{4 cores were
dedicated to the Host}, and we prevented tasks to be scheduled on other
cores by adopting some of the OS-level isolation techniques we already
used: Tickless Mode, RCU callback offloading and IRQ affinity. Kernel
scheduling isolation was not used because it interfered with the way KVM
assigned virtual CPUs to phisical ones. The analyzed variables concern
the introduction of a stress load (noise) of varying size originating
from another Guest VM and the application of vCPU pinning (crucial for
avoiding context migrations and stabilizing latencies).

To ensure strict real-time conditions and accurate latency measurements,
the execution procedure was carefully standardized across all test
scenarios. The process involved specific compilation flags, scheduler
modifications, and rigid execution parameters designed to eliminate
typical operating system interference.

\textbf{Compilation and Preparation} Each of the benchmarks were
modified to support automatically changing scheduling priority and
repeated executions (the source code is available in this repository).
We then compiled them using \texttt{-O2} optimizations and linked with
the necessary POSIX real-time and threading libraries:

\begin{Shaded}
\begin{Highlighting}[]
\FunctionTok{gcc}\NormalTok{ {-}O2 {-}o test3 test3.c {-}lpthread {-}lrt}
\end{Highlighting}
\end{Shaded}

\textbf{Benchmark Execution Parameters} The TACLe benchmarks were
executed via the command line with a strict set of arguments to maintain
deterministic behavior. The fundamental flags applied across the tests
included:

\begin{itemize}
\item
  \texttt{-\/-mlockall}: Locks the process's memory pages directly into
  RAM, preventing any unpredictable latency spikes caused by page faults
  or disk swapping.
\item
  \texttt{-\/-priority=99}: Enforces the highest real-time priority
  internally for the benchmark's execution thread.
\item
  \texttt{-\/-affinity=1}: Applies CPU pinning, locking the execution
  strictly to CPU core 1. This is a critical step to prevent costly
  context migrations across different processor cores.
\end{itemize}

Below are the exact execution commands utilized for the targeted
benchmarks:

\begin{Shaded}
\begin{Highlighting}[]
\CommentTok{\# Executing test3 (pseudo{-}cyclictest) for 10,000 loops and outputting a latency histogram}
\FunctionTok{sudo}\NormalTok{ ./test3 {-}{-}mlockall {-}{-}priority=99 {-}{-}affinity=1 {-}{-}loops=10000 {-}{-}histogram=1000000 {-}{-}histfile=results/results\_test3\_baseline.log}
\end{Highlighting}
\end{Shaded}

\textbf{Noise Generation Strategy} To evaluate the architectural
robustness and jitter expansion during the noisy scenarios, interference
was artificially injected into the system. This was achieved using
\texttt{stress-ng} to spawn multiple aggressive workers, intentionally
taxing the CPU cores and the virtual memory subsystem to simulate severe
cache thrashing and scheduling pressure:

\begin{Shaded}
\begin{Highlighting}[]
\ExtensionTok{stress{-}ng}\NormalTok{ {-}{-}cpu 16 {-}{-}vm 16 {-}{-}vm{-}bytes 1G {-}{-}timeout 10m}
\end{Highlighting}
\end{Shaded}

\textbf{Virtual Machine Configuration} To clearly understand the extent
of the interference of a different Guest VM stress workload on the
critical guest, we designed two test scenarios:

\begin{itemize}
\item
  \textbf{Small Noise}: a small Guest VM with just 2 vCPUs and 4 GB of
  RAM. This would introduce a small amount of noise, thus testing if at
  least some isolation is provided.
\item
  \textbf{Big Noise}: a more beefy Guest VM configuration with 16 vCPUs
  and 16 GB of RAM. In this case, the interference introduced is more
  severe, pushing the limits of the isolation mechanism.
\item
  \textbf{Big Noise with pinning}: the same of the previous setup but
  with added static vCPU pinning. This would determine how much the
  scheduling algorithm affects the isolation of the virtual machines.
\end{itemize}

\subsubsection{TACLeBench Baseline Execution
Analysis}\label{taclebench-baseline-execution-analysis}

Establishing this baseline is critical for evaluating the determinism
and worst-case execution time (WCET) latencies of the KVM hypervisor.
The following analysis interprets the raw execution logs to quantify the
baseline scheduling determinism before the introduction of concurrent
noisy guests.

\paragraph{Baseline Latency Summary}\label{baseline-latency-summary}

\begin{table*}
\centering
\caption{KVM Baseline Latency Summary}
\label{tab:kvm:baseline}
\begin{tabularx}{\linewidth}{@{}YYYYYY@{}}
\toprule
Benchmark & Total Loops & Min Latency ($\mu s$) & Avg Latency ($\mu s$) & Max Latency ($\mu s$) & Absolute Jitter (Max - Min)\tabularnewline
\midrule
\textbf{DEBIE} & 10,000 & 23,578 & 27,214 & 32,409 & 8,831 $\mu s$\tabularnewline
\textbf{huffenc} & 1,000,000 & 15 & 17 & 65 & 50 $\mu s$\tabularnewline
\textbf{lift} & 1,000,000 & 19 & 20 & 57 & 38 $\mu s$\tabularnewline
\textbf{matrix1} & 1,000,000 & 0 & 0 & 1 & 1 $\mu s$\tabularnewline
\textbf{test3} & 10,000 & 8,169 & 8,363 & 9,471 & 1,302 $\mu s$\tabularnewline
\bottomrule
\end{tabularx}
\end{table*}

\paragraph{Workload-Specific
Behavior}\label{workload-specific-behavior}

\textbf{DEBIE}

\begin{itemize}
\item
  The DEBIE benchmark represents the heaviest parallel workload,
  executing 10,000 loops with execution times ranging from 23,578 $\mu s$ to
  32,409 $\mu s$.
\item
  The absolute jitter sits at 8,831 $\mu s$, indicating typical baseline
  scheduling overhead and preemption for a heavy multi-tasking
  application on an unisolated KVM host.
\end{itemize}

\textbf{Huffenc}

\begin{itemize}
\item
  The \texttt{huffenc} benchmark executed 1,000,000 loops with a stable
  average latency of 17 $\mu s$.
\item
  The maximum latency peaked at 65 $\mu s$, resulting in an absolute jitter
  of 50 $\mu s$.
\end{itemize}

\textbf{Lift}

\begin{itemize}
\item
  The \texttt{lift} application benchmark showed strong determinism,
  maintaining an average latency of 20 $\mu s$ close to its minimum latency.
\item
  The worst-case latency was 57 $\mu s$, yielding an absolute jitter of 38 $\mu s$
  on this KVM baseline.
\end{itemize}

\textbf{Matrix1}

\begin{itemize}
\item
  The \texttt{matrix1} kernel benchmark was extremely fast, executing
  1,000,000 loops with sub-microsecond average latency.
\item
  The maximum execution time was strictly capped at just 1 $\mu s$,
  demonstrating near-perfect cache efficiency and virtually zero jitter
  (1 $\mu s$).
\end{itemize}

\textbf{Test3}

\begin{itemize}
\item
  The \texttt{test3} benchmark executed 10,000 loops with an average
  latency of 8,363 $\mu s$.
\item
  The maximum latency reached 9,471 $\mu s$, resulting in an absolute jitter
  of 1,302 $\mu s$, representing standard noise interference on a
  medium-weight computational loop.
\end{itemize}

\paragraph{Real-Time Systems
Assessment}\label{real-time-systems-assessment}

This KVM baseline provides the reference point for unisolated
virtualized performance. Lightweight tasks (\texttt{matrix1}) exhibit
virtually zero jitter (1 $\mu s$ peak), \texttt{lift} shows a moderate
baseline jitter of 38 $\mu s$ s, while medium and heavy tasks (\texttt{test3},
\texttt{DEBIE}) show the natural variance introduced by the standard KVM
scheduler (up to 8,831 $\mu s$ s of jitter for DEBIE). These metrics will be
crucial for quantifying the exact determinism improvements when CPU
pinning and LL-RT kernel isolation techniques are introduced.

\subsubsection{TACLeBench Small Noise Execution
Analysis}\label{taclebench-small-noise-execution-analysis}

This analysis evaluates the determinism and worst-case execution time
(WCET) latencies of the KVM hypervisor under a ``Small Noise''
configuration. By introducing minor background system stress, we can
quantify the sensitivity of the unisolated KVM scheduler to light
interference.

\paragraph{Small Noise Latency
Summary}\label{small-noise-latency-summary}

\begin{table*}
\centering
\caption{KVM Small Noise Latency Summary}
\label{tab:kvm:small-noise}
\begin{tabularx}{\linewidth}{@{}YYYYYY@{}}
\toprule
Benchmark & Total Loops & Min Latency ($\mu s$) & Avg Latency ($\mu s$) & Max Latency ($\mu s$) & Absolute Jitter (Max - Min)\tabularnewline
\midrule
\textbf{DEBIE} & 10,000 & 23,800 & 27,221 & 30,872 & 7,072 $\mu s$\tabularnewline
\textbf{huffenc} & 1,000,000 & 15 & 17 & 49 & 34 $\mu s$\tabularnewline
\textbf{lift} & 1,000,000 & 19 & 20 & 55 & 36 $\mu s$\tabularnewline
\textbf{matrix1} & 1,000,000 & 0 & 0 & 1 & 1 $\mu s$\tabularnewline
\textbf{test3} & 10,000 & 8,078 & 8,270 & 8,612 & 534 $\mu s$\tabularnewline
\bottomrule
\end{tabularx}
\end{table*}

\paragraph{Workload-Specific
Behavior}\label{workload-specific-behavior-1}

\textbf{DEBIE}

\begin{itemize}
\item
  The heavy parallel DEBIE workload experiences execution times ranging
  from 23,800 $\mu s$ to 30,872 $\mu s$.
\item
  The absolute jitter sits at 7,072 $\mu s$, showing that even small amounts
  of noise can cause significant preemption and scheduling variation for
  complex multi-tasking applications.
\end{itemize}

\textbf{Huffenc}

\begin{itemize}
\item
  The sequential \texttt{huffenc} task executes with an average latency
  of 17 $\mu s$ and a maximum peak of 49 $\mu s$.
\item
  This results in an absolute jitter of 34 $\mu s$.
\end{itemize}

\textbf{Lift}

\begin{itemize}
\item
  The \texttt{lift} control benchmark maintains strong determinism under
  light noise, averaging 20 $\mu s$ with a worst-case execution time of 55
  $\mu s$.
\item
  The resulting jitter is bounded at 36 $\mu s$, showing good resilience to
  minor disturbances, and is close to its already-moderate baseline
  jitter (38 $\mu s$).
\end{itemize}

\textbf{Matrix1}

\begin{itemize}
\item
  The lightweight, cache-bound \texttt{matrix1} benchmark performs
  almost flawlessly under the small noise configuration, logging
  sub-microsecond average execution time with a max of just 1 $\mu s$.
\item
  This yields a minimal absolute jitter of 1 $\mu s$, indicating negligible
  cache thrashing or scheduling interruptions occurred during its
  1,000,000 loops.
\end{itemize}

\textbf{Test3}

\begin{itemize}
\item
  The \texttt{test3} WCET stress test executes with an average latency
  of 8,270 $\mu s$.
\item
  The maximum latency reaches 8,612 $\mu s$, keeping the absolute jitter
  contained at just 534 $\mu s$.
\end{itemize}

\paragraph{Real-Time Systems
Assessment}\label{real-time-systems-assessment-1}

Under the KVM ``Small Noise'' scenario, the hypervisor exhibits
relatively stable behavior for lightweight and medium workloads.
Notably, \texttt{matrix1} shows near-perfect execution with negligible
jitter (1 $\mu s$ peak), and \texttt{test3} maintains tight WCET margins.
However, the \texttt{DEBIE} workload still suffers from over 7
milliseconds of jitter, confirming that unisolated schedulers struggle
to guarantee execution determinism for long-running, parallel tasks even
when background noise is minimal.

\subsubsection{TACLeBench Big Noise Execution
Analysis}\label{taclebench-big-noise-execution-analysis}

This analysis evaluates the determinism and worst-case execution time
(WCET) latencies of the KVM hypervisor under a ``Big Noise''
configuration. By introducing heavy simulated stress via
\texttt{stress-ng}, we can quantify how severely external disturbances
and resource contention degrade the scheduling predictability of the
baseline KVM setup before any CPU pinning or kernel isolation is
applied.

\paragraph{Big Noise Latency Summary}\label{big-noise-latency-summary}

\begin{table*}
\centering
\caption{KVM Big Noise Latency Summary}
\label{tab:kvm:big-noise}
\begin{tabularx}{\linewidth}{@{}YYYYYY@{}}
\toprule
Benchmark & Total Loops & Min Latency ($\mu s$) & Avg Latency ($\mu s$) & Max Latency ($\mu s$) & Absolute Jitter (Max - Min)\tabularnewline
\midrule
\textbf{DEBIE} & 10,000 & 24,773 & 28,845 & 43,849 & 19,076 $\mu s$\tabularnewline
\textbf{huffenc} & 1,000,000 & 15 & 17 & 393 & 378 $\mu s$\tabularnewline
\textbf{lift} & 1,000,000 & 19 & 21 & 186 & 167 $\mu s$\tabularnewline
\textbf{matrix1} & 1,000,000 & 0 & 0 & 1 & 1 $\mu s$\tabularnewline
\textbf{test3} & 10,000 & 8,243 & 8,457 & 9,458 & 1,215 $\mu s$\tabularnewline
\bottomrule
\end{tabularx}
\end{table*}

\paragraph{Workload-Specific
Behavior}\label{workload-specific-behavior-2}

\textbf{DEBIE}

\begin{itemize}
\item
  The heavy parallel DEBIE workload experiences significant disruption
  under noise, with the maximum execution time expanding to 43,849 $\mu s$.
\item
  The absolute jitter nearly doubles compared to the KVM baseline,
  reaching 19,076 $\mu s$, highlighting severe preemption and scheduling
  delays when competing for resources.
\end{itemize}

\textbf{Huffenc}

\begin{itemize}
\item
  While the minimum latency for the \texttt{huffenc} task remains stable
  at 15 $\mu s$, the average latency stays close to baseline at 17 $\mu s$ .
\item
  The maximum latency, however, spikes dramatically to 393 $\mu s$, resulting
  in an absolute jitter of 378 $\mu s$ --- the largest percentage increase
  (+501\% on WCET) of any benchmark tested, revealing severe worst-case
  tail latency under heavy stress despite a stable average.
\end{itemize}

\textbf{Lift}

\begin{itemize}
\item
  The \texttt{lift} benchmark shows a sharp breakdown in worst-case
  determinism, with the maximum latency spiking to 186 $\mu s$.
\item
  This generates an absolute jitter of 167 $\mu s$ (+226\% on WCET versus
  baseline), indicating that even highly repetitive, lightweight control
  tasks suffer from severe, albeit occasional, micro-interruptions under
  heavy KVM load.
\end{itemize}

\textbf{Matrix1}

\begin{itemize}
\item
  The \texttt{matrix1} execution remains extremely fast, maintaining a
  sub-microsecond average latency.
\item
  The maximum latency stays at 1 $\mu s$, representing a jitter of just 1 $\mu s$
  --- a modest but measurable +32\% increase in WCET versus baseline,
  indicating minor cache thrashing or context-switching interference
  affecting the otherwise near-instantaneous execution.
\end{itemize}

\textbf{Test3}

\begin{itemize}
\item
  The \texttt{test3} benchmark registers an average latency of 8,457 $\mu s$
  and a worst-case peak of 9,458 $\mu s$.
\item
  The absolute jitter sits at 1,215 $\mu s$, which is surprisingly consistent
  with its baseline performance, suggesting this specific computational
  loop is somewhat resilient to the injected memory and CPU stress.
\end{itemize}

\paragraph{Real-Time Systems
Assessment}\label{real-time-systems-assessment-2}

The KVM ``Big Noise'' scenario illustrates the vulnerability of
unisolated hypervisor scheduling. While \texttt{test3} and
\texttt{matrix1} show relative resilience, the heavy \texttt{DEBIE}
workload suffers massive jitter expansion (19,076 $\mu s$), and the
lightweight \texttt{huffenc} and \texttt{lift} control loops experience
the sharpest \emph{percentage} degradation of the whole suite:
\texttt{huffenc} WCET jumps from 65 $\mu s$ to 393 $\mu s$ (+501\%), and
\texttt{lift} WCET jumps from 57 $\mu s$ to 186 $\mu s$ (+226\%), despite both
keeping stable average latencies. This shows that unisolated KVM
scheduling can produce severe worst-case tail latency spikes on short,
high-frequency workloads even when average-case behavior looks stable
--- a critical finding for real-time guarantees. These results suggest
the lack of boundaries required for real-time applications.

\subsubsection{TACLeBench Big Noise Pinned Execution
Analysis}\label{taclebench-big-noise-pinned-execution-analysis}

This phase of the evaluation investigates the determinism and worst-case
execution time (WCET) latencies of the KVM hypervisor under a ``Big
Noise Pinned'' configuration. Because the unpinned ``Big Noise'' results
exhibited high latencies and severe jitter, CPU pinning was introduced
to lock the execution to specific cores. The following analysis examines
the execution logs to quantify how this pinning mitigates system
disturbances under KVM.

\paragraph{Big Noise Pinned Latency
Summary}\label{big-noise-pinned-latency-summary}

\begin{table*}
\centering
\caption{KVM Big Noise Pinned Latency Summary}
\label{tab:kvm:big-noise-pinned}
\begin{tabularx}{\linewidth}{@{}YYYYYY@{}}
\toprule
Benchmark & Total Loops & Min Latency ($\mu s$) & Avg Latency ($\mu s$) & Max Latency ($\mu s$) & Absolute Jitter (Max - Min)\tabularnewline
\midrule
\textbf{DEBIE} & 10,000 & 23,893 & 27,444 & 37,284 & 13,391 $\mu s$\tabularnewline
\textbf{huffenc} & 1,000,000 & 15 & 17 & 73 & 58 $\mu s$\tabularnewline
\textbf{lift} & 1,000,000 & 19 & 21 & 70 & 51 $\mu s$\tabularnewline
\textbf{matrix1} & 1,000,000 & 0 & 0 & 1 & 1 $\mu s$\tabularnewline
\textbf{test3} & 10,000 & 8,544 & 8,745 & 10,996 & 2,452 $\mu s$\tabularnewline
\bottomrule
\end{tabularx}
\end{table*}

\paragraph{Workload-Specific
Behavior}\label{workload-specific-behavior-3}

\textbf{DEBIE}

\begin{itemize}
\item
  The DEBIE benchmark ranges from a minimum of 23,893 $\mu s$ to a maximum of
  37,284 $\mu s$.
\item
  Pinning the execution restrains the absolute jitter to 13,391 $\mu s$.
  While still representing a significant preemption delay, it is a
  notable improvement over the KVM unpinned big noise scenario.
\end{itemize}

\textbf{Huffenc}

\begin{itemize}
\item
  The \texttt{huffenc} task executes with an average latency of 16 $\mu s$
  and a maximum peak of 62 $\mu s$.
\item
  CPU pinning dramatically reduces the absolute jitter to 58 $\mu s$, and ---
  critically --- pulls the WCET down from 393 $\mu s$ (unpinned Big Noise)
  back to 73 $\mu s$, restoring a high degree of stability to this sequential
  loop.
\end{itemize}

\textbf{Lift}

\begin{itemize}
\item
  Under the pinned configuration, the \texttt{lift} benchmark regains
  much of its determinism, maintaining an average latency of 21 $\mu s$.
\item
  The maximum latency reaches 70 $\mu s$, yielding an absolute jitter of 51
  $\mu s$ --- far below the 167 $\mu s$ jitter seen in the unpinned Big Noise
  test, and now close to the baseline's own jitter of 38 $\mu s$.
\end{itemize}

\textbf{Matrix1}

\begin{itemize}
\item
  The \texttt{matrix1} execution is almost perfectly insulated by the
  pinning, maintaining a sub-microsecond average latency and an absolute
  worst-case execution time of just 1 $\mu s$.
\item
  This 1 $\mu s$ absolute jitter confirms that cache thrashing and context
  migrations have been largely mitigated, in line with its baseline
  behavior.
\end{itemize}

\textbf{Test3}

\begin{itemize}
\item
  The \texttt{test3} benchmark registers an average latency of 8,745 $\mu s$
  and a worst-case peak of 10,996 $\mu s$.
\item
  With an absolute jitter of 2,452 $\mu s$, this computational loop maintains
  a relatively controlled variance compared to unpinned scenarios.
\end{itemize}

\paragraph{Real-Time Systems
Assessment}\label{real-time-systems-assessment-3}

The KVM ``Big Noise Pinned'' baseline demonstrates the critical
importance of CPU pinning when operating in highly congested
environments. By binding tasks to specific cores, the KVM scheduler
prevents the catastrophic latency spikes observed in the unpinned tests
--- most notably for \texttt{huffenc}, whose WCET drops from 393 $\mu s$ back
to 73 $\mu s$ once pinning is applied. Lightweight tasks (\texttt{matrix1},
\texttt{huffenc}, \texttt{lift}) return to near-baseline determinism,
with \texttt{lift}'s WCET (70 $\mu s$) settling close to its own baseline (57
$\mu s$) rather than the elevated 186 $\mu s$ seen unpinned. Heavier workloads
(\texttt{DEBIE}) see their jitter margins compressed significantly. This
confirms that static pinning is a highly effective first step in
isolating real-time workloads on KVM before applying further
kernel-level techniques.

\subsubsection{Max Latency Summary
($\mu s$)}\label{max-latency-summary-ux3bcs}

To quickly evaluate system stability, the following table exclusively
reports the peak values (\textbf{Max Latency}). This format allows for
an at-a-glance comparison of the Worst-Case Execution Time across the
four operational scenarios, directly highlighting the impact of noise
and the effectiveness of CPU pinning in containing interference.

\begin{table*}
\centering
\caption{KVM WCET Peak Values Summary}
\label{tab:kvm:wcet-peak}
\begin{tabularx}{\linewidth}{@{}YYYYY@{}}
\toprule
Benchmark & Baseline & Small Noise & Big Noise & Big Noise Pinned\tabularnewline
\midrule
\textbf{DEBIE} & 32,409 & 30,872 & 43,849 & 37,284\tabularnewline
\textbf{huffenc} & 65 & 49 & 393 & 73\tabularnewline
\textbf{lift} & 57 & 55 & 186 & 70\tabularnewline
\textbf{matrix1} & 1 & 1 & 1 & 1\tabularnewline
\textbf{test3} & 9,471 & 8,612 & 9,458 & 10,996\tabularnewline
\bottomrule
\end{tabularx}
\end{table*}

Furthermore, alongside the WCET summary, this section now includes a
detailed breakdown of the percentage increments for both the average
latency (expressed as Mean \(\pm\) Standard Deviation) and the maximum
latency (WCET). This addition provides a precise quantitative analysis
of the performance degradation induced by the heavy background noise
compared to the baseline execution for each individual benchmark.

\paragraph{DEBIE}\label{debie}

\begin{table*}
\centering
\caption{KVM DEBIE Execution Time Statistics}
\label{tab:kvm:debie}
\begin{tabularx}{\linewidth}{@{}YYYYY@{}}
\toprule
METRIC & BASELINE & SMALL NOISE & BIG NOISE & BIG NOISE PINNED\tabularnewline
\midrule
\textbf{Mean ± SD} & 27214.36 ± 1303.61 $\mu s$ & 27221.56 ± 1304.61 $\mu s$ & 28845.81 ± 3813.54 $\mu s$ & 27444.09 ± 1327.72 $\mu s$\tabularnewline
\textbf{WCET (Max)} & 32409 $\mu s$ & 30872 $\mu s$ & 43849 $\mu s$ & 37284 $\mu s$\tabularnewline
\bottomrule
\end{tabularx}
\end{table*}

\textbf{PERCENTAGE INCREMENTS (Stress vs.~Baseline):} * Average
Increment: +5.99\% * WCET Increment: +35.30\%

\begin{figure}
\centering
\includesvg[width=\linewidth]{../KVM/tests_TACLe/plots/kvm_TACLe_debie_boxplot.svg}
\caption{TACLe benchmark - debie execution time on KVM}
\end{figure}

\paragraph{HUFFENC}\label{huffenc}

\begin{table*}
\centering
\caption{KVM huff\_enc Execution Time Statistics}
\label{tab:kvm:huffenc}
\begin{tabularx}{\linewidth}{@{}YYYYY@{}}
\toprule
METRIC & BASELINE & SMALL NOISE & BIG NOISE & BIG NOISE PINNED\tabularnewline
\midrule
\textbf{Mean ± SD} & 16619.81 ± 3708.79 ns & 17306.98 ± 4400.23 ns & 16624.73 ± 3207.31 ns & 16871.85 ± 3772.31 ns\tabularnewline
\textbf{WCET (Max)} & 65431 ns & 393480 ns & 49307 ns & 72590 ns\tabularnewline
\bottomrule
\end{tabularx}
\end{table*}

\textbf{PERCENTAGE INCREMENTS (Stress vs.~Baseline):} * Average
Increment: +4.13\% * WCET Increment: +501.37\%

\begin{figure}
\centering
\includesvg[width=\linewidth]{../KVM/tests_TACLe/plots_nanosec/svg/kvm_TACLe_huffenc_boxplot.svg}
\caption{TACLe benchmark - huff\_enc execution time on KVM}
\end{figure}

\paragraph{LIFT}\label{lift}

\begin{table*}
\centering
\caption{KVM lift Execution Time Statistics}
\label{tab:kvm:lift}
\begin{tabularx}{\linewidth}{@{}YYYYY@{}}
\toprule
METRIC & BASELINE & SMALL NOISE & BIG NOISE & BIG NOISE PINNED\tabularnewline
\midrule
\textbf{Mean ± SD} & 20106.37 ± 2997.02 ns & 20113.88 ± 3034.06 ns & 21218.61 ± 4413.84 ns & 21259.27 ± 5273.76 ns\tabularnewline
\textbf{WCET (Max)} & 57169 ns & 54767 ns & 186424 ns & 69974 ns\tabularnewline
\bottomrule
\end{tabularx}
\end{table*}

\textbf{PERCENTAGE INCREMENTS (Stress vs.~Baseline):} * Average
Increment: +5.53\% * WCET Increment: +226.09\%

\begin{figure}
\centering
\includesvg[width=\linewidth]{../KVM/tests_TACLe/plots_nanosec/svg/kvm_TACLe_lift_boxplot.svg}
\caption{TACLe benchmark - lift execution time on KVM}
\end{figure}

\paragraph{MATRIX1}\label{matrix1}

\begin{table*}
\centering
\caption{KVM matrix1 Execution Time Statistics}
\label{tab:kvm:matrix1}
\begin{tabularx}{\linewidth}{@{}YYYYY@{}}
\toprule
METRIC & BASELINE & SMALL NOISE & BIG NOISE & BIG NOISE PINNED\tabularnewline
\midrule
\textbf{Mean ± SD} & 367.88 ± 6.57 ns & 369.95 ± 6.66 ns & 370.74 ± 20.92 ns & 384.14 ± 29.37 ns\tabularnewline
\textbf{WCET (Max)} & 720 ns & 790 ns & 950 ns & 1070 ns\tabularnewline
\bottomrule
\end{tabularx}
\end{table*}

\textbf{PERCENTAGE INCREMENTS (Stress vs.~Baseline):} * Average
Increment: +0.78\% * WCET Increment: +31.94\%

\begin{figure}
\centering
\includesvg[width=\linewidth]{../KVM/tests_TACLe/plots_nanosec/svg/kvm_TACLe_matrix1_boxplot.svg}
\caption{TACLe benchmark - matrix1 execution time on KVM}
\end{figure}

\paragraph{TEST3}\label{test3}

\begin{table*}
\centering
\caption{KVM test3 Execution Time Statistics}
\label{tab:kvm:test3}
\begin{tabularx}{\linewidth}{@{}YYYYY@{}}
\toprule
METRIC & BASELINE & SMALL NOISE & BIG NOISE & BIG NOISE PINNED\tabularnewline
\midrule
\textbf{Mean ± SD} & 8363.95 ± 31.59 $\mu s$ & 8270.33 ± 24.90 $\mu s$ & 8457.93 ± 52.13 $\mu s$ & 8745.21 ± 87.11 $\mu s$\tabularnewline
\textbf{WCET (Max)} & 9471 $\mu s$ & 8612 $\mu s$ & 9458 $\mu s$ & 10996 $\mu s$\tabularnewline
\bottomrule
\end{tabularx}
\end{table*}

\textbf{PERCENTAGE INCREMENTS (Stress vs.~Baseline):} * Average
Increment: +1.12\% * WCET Increment: -0.14\%

\begin{figure}
\centering
\includesvg[width=\linewidth]{../KVM/tests_TACLe/plots/kvm_TACLe_test3_boxplot.svg}
\caption{TACLe benchmark - test3 execution time on KVM}
\end{figure}
